% GENERATED FILE. Edit canonical lessons, then run npm run generate:books.
% canonical-source-hash: fnv1a64:34bba5c0f42bb693
% canonical-lessons: MR-C523-galica, MR-C523-rajai, MR-C523-macchardani, MR-C523-catai, MR-C523-gundalne

\chapter{Carpets and Quilts}
\label{ch:mr-carpets-and-quilts}
\hlchaptermodality{523}

\begin{chapteropening}
\textbf{By the end of this chapter:} \emph{I can talk about carpets, quilts, nets and mats.}

Complete the last lesson of chapter 523: I can talk about carpets, quilts, nets and mats.
\end{chapteropening}

\section[\texorpdfstring{\mr{गालिचा} (gālicā)}{गालिचा (gālicā)}]{\mr{गालिचा} (gālicā) — a carpet; a rug}

\label{lesson:MR-C523-galica}



\begin{quote}
Before the new one: say the Marathi for a floor tile, then the Marathi for to hang up.
\end{quote}

\subsection*{You'll want to know: \mr{गालिचा}}
\textbf{\mr{गालिचा}} — \emph{gālicā} — \textquotedblleft{}a carpet; a rug\textquotedblright{}.

\begin{quote}
\mr{दिवाणखान्यात} \mr{लाल} \mr{गालिचा} \mr{आहे}.
\end{quote}

\begin{quote}
\emph{divāṇkhānyāt lāl gālicā āhe.}
\end{quote}

\begin{quote}
There is a red carpet in the living room.
\end{quote}

Masculine.

\begin{cousinweb}
\textbf{\mr{गालिचा}} (\emph{gālicā}) is a loan from Persian \emph{qālīča}, \textquotedblleft{}a small carpet\textquotedblright{}.
\end{cousinweb}

\begin{grammarlens}[title={What you've built}]
One more word for describing what you lived and felt.
\end{grammarlens}

\subsection*{Your turn}
\begin{itemize}
\raggedright
  \item \emph{Say it:} \emph{gālicā}
  \item \emph{Say it:} \emph{gālicā}, once more
  \item \emph{Say it:} \emph{divāṇkhānyāt lāl gālicā āhe.} — There is a red carpet in the living room.
  \item \emph{Recall:} say the Marathi for a floor tile, then the Marathi for to hang up, then say \emph{gālicā} again
\end{itemize}

\begin{tcolorbox}[breakable,colback=teal!4,colframe=teal!35!black,title={Before you move on}]
What does \mr{गालिचा} mean? (\textquotedblleft{}A carpet; a rug\textquotedblright{}.) Say it once more.
\end{tcolorbox}
\section[\texorpdfstring{\mr{रजई} (rajaī)}{रजई (rajaī)}]{\mr{रजई} (rajaī) — a quilt}

\label{lesson:MR-C523-rajai}



\begin{quote}
Before the new one: say the Marathi for to hang up, then the Marathi for a carpet.
\end{quote}

\subsection*{You'll want to know: \mr{रजई}}
\textbf{\mr{रजई}} — \emph{rajaī} — \textquotedblleft{}a quilt\textquotedblright{}.

\begin{quote}
\mr{थंडीत} \mr{आजी} \mr{जाड} \mr{रजई} \mr{घेते}.
\end{quote}

\begin{quote}
\emph{thaṇḍīt ājī jāḍ rajaī ghete.}
\end{quote}

\begin{quote}
In the cold, grandmother covers herself with a thick quilt.
\end{quote}

Feminine.

\begin{cousinweb}
\textbf{\mr{रजई}} (\emph{rajaī}) is usually traced to Persian \emph{razāī}, \textquotedblleft{}a quilt\textquotedblright{}.
\end{cousinweb}

\begin{grammarlens}[title={What you've built}]
One more word for describing what you lived and felt.
\end{grammarlens}

\subsection*{Your turn}
\begin{itemize}
\raggedright
  \item \emph{Say it:} \emph{rajaī}
  \item \emph{Say it:} \emph{rajaī}, once more
  \item \emph{Say it:} \emph{thaṇḍīt ājī jāḍ rajaī ghete.} — In the cold, grandmother covers herself with a thick quilt.
  \item \emph{Recall:} say the Marathi for to hang up, then the Marathi for a carpet, then say \emph{rajaī} again
\end{itemize}

\begin{tcolorbox}[breakable,colback=teal!4,colframe=teal!35!black,title={Before you move on}]
What does \mr{रजई} mean? (\textquotedblleft{}A quilt\textquotedblright{}.) Say it once more.
\end{tcolorbox}
\section[\texorpdfstring{\mr{मच्छरदाणी} (macchardāṇī)}{मच्छरदाणी (macchardāṇī)}]{\mr{मच्छरदाणी} (macchardāṇī) — a mosquito net}

\label{lesson:MR-C523-macchardani}



\begin{quote}
Before the new one: say the Marathi for a carpet, then the Marathi for a quilt.
\end{quote}

\subsection*{You'll want to know: \mr{मच्छरदाणी}}
\textbf{\mr{मच्छरदाणी}} — \emph{macchardāṇī} — \textquotedblleft{}a mosquito net\textquotedblright{}.

\begin{quote}
\mr{डास} \mr{खूप} \mr{होते}, \mr{म्हणून} \mr{आम्ही} \mr{मच्छरदाणी} \mr{लावली}.
\end{quote}

\begin{quote}
\emph{ḍās khūp hote, mhaṇūn āmhī macchardāṇī lāvlī.}
\end{quote}

\begin{quote}
There were lots of mosquitoes, so we put up a mosquito net.
\end{quote}

Feminine.

\mr{मच्छरदाणी} \mr{लावणे}, to put up a net.

\begin{cousinweb}
\textbf{\mr{मच्छरदाणी}} (\emph{macchardāṇī}) is a compound of \mr{मच्छर}, \textquotedblleft{}mosquito\textquotedblright{}, and -\mr{दाणी}, \textquotedblleft{}a holder\textquotedblright{}; the origin of \mr{मच्छर} is not settled.
\end{cousinweb}

\begin{grammarlens}[title={What you've built}]
One more word for describing what you lived and felt.
\end{grammarlens}

\subsection*{Your turn}
\begin{itemize}
\raggedright
  \item \emph{Say it:} \emph{macchardāṇī}
  \item \emph{Say it:} \emph{macchardāṇī}, once more
  \item \emph{Say it:} \emph{ḍās khūp hote, mhaṇūn āmhī macchardāṇī lāvlī.} — There were lots of mosquitoes, so we put up a mosquito net.
  \item \emph{Recall:} say the Marathi for a carpet, then the Marathi for a quilt, then say \emph{macchardāṇī} again
\end{itemize}

\begin{tcolorbox}[breakable,colback=teal!4,colframe=teal!35!black,title={Before you move on}]
What does \mr{मच्छरदाणी} mean? (\textquotedblleft{}A mosquito net\textquotedblright{}.) Say it once more.
\end{tcolorbox}
\section[\texorpdfstring{\mr{चटई} (caṭaī)}{चटई (caṭaī)}]{\mr{चटई} (caṭaī) — a mat}

\label{lesson:MR-C523-catai}



\begin{quote}
Before the new one: say the Marathi for a quilt, then the Marathi for a mosquito net.
\end{quote}

\subsection*{You'll want to know: \mr{चटई}}
\textbf{\mr{चटई}} — \emph{caṭaī} — \textquotedblleft{}a mat\textquotedblright{}.

\begin{quote}
\mr{मुलं} \mr{चटईवर} \mr{बसून} \mr{जेवली}.
\end{quote}

\begin{quote}
\emph{mulaṁ caṭaīvar basūn jevlī.}
\end{quote}

\begin{quote}
The children sat on the mat and ate.
\end{quote}

Feminine.

\begin{cousinweb}
\textbf{\mr{चटई}} (\emph{caṭaī}) is inherited: CDIAL lists Marathi \emph{caṭāī}, \textquotedblleft{}matting\textquotedblright{}, under \emph{*caṭṭa}, \textquotedblleft{}bamboo matting\textquotedblright{} (CDIAL 4574).
\end{cousinweb}

\begin{grammarlens}[title={What you've built}]
One more word for describing what you lived and felt.
\end{grammarlens}

\subsection*{Your turn}
\begin{itemize}
\raggedright
  \item \emph{Say it:} \emph{caṭaī}
  \item \emph{Say it:} \emph{caṭaī}, once more
  \item \emph{Say it:} \emph{mulaṁ caṭaīvar basūn jevlī.} — The children sat on the mat and ate.
  \item \emph{Recall:} say the Marathi for a quilt, then the Marathi for a mosquito net, then say \emph{caṭaī} again
\end{itemize}

\begin{tcolorbox}[breakable,colback=teal!4,colframe=teal!35!black,title={Before you move on}]
What does \mr{चटई} mean? (\textquotedblleft{}A mat\textquotedblright{}.) Say it once more.
\end{tcolorbox}
\section[\texorpdfstring{\mr{गुंडाळणे} (guṇḍāḷṇe)}{गुंडाळणे (guṇḍāḷṇe)}]{\mr{गुंडाळणे} (guṇḍāḷṇe) — to roll up; to wrap}

\label{lesson:MR-C523-gundalne}



\begin{quote}
Before the new one: say the Marathi for a mosquito net, then the Marathi for a mat.
\end{quote}

\subsection*{You'll want to know: \mr{गुंडाळणे}}
\textbf{\mr{गुंडाळणे}} — \emph{guṇḍāḷṇe} — \textquotedblleft{}to roll up; to wrap\textquotedblright{}.

\begin{quote}
\mr{सकाळी} \mr{त्याने} \mr{गालिचा} \mr{गुंडाळला} \mr{आणि} \mr{बाहेर} \mr{नेला}.
\end{quote}

\begin{quote}
\emph{sakāḷī tyāne gālicā guṇḍāḷlā āṇi bāher nelā.}
\end{quote}

\begin{quote}
In the morning he rolled up the carpet and took it outside.
\end{quote}

Past, with -\mr{ने} (he …): \mr{त्याने} \mr{चटई} \mr{गुंडाळली}, \emph{tyāne caṭaī guṇḍāḷlī}.

\begin{cousinweb}
\textbf{\mr{गुंडाळणे}} (\emph{guṇḍāḷṇe}) is a Marathi verb for rolling or wrapping; its origin is not settled.
\end{cousinweb}

\begin{grammarlens}[title={What you've built}]
One more word for describing what you lived and felt.
\end{grammarlens}

\subsection*{Your turn}
\begin{itemize}
\raggedright
  \item \emph{Say it:} \emph{guṇḍāḷṇe}
  \item \emph{Say it:} \emph{guṇḍāḷṇe}, once more
  \item \emph{Say it:} \emph{sakāḷī tyāne gālicā guṇḍāḷlā āṇi bāher nelā.} — In the morning he rolled up the carpet and took it outside.
  \item \emph{Recall:} say the Marathi for a mosquito net, then the Marathi for a mat, then say \emph{guṇḍāḷṇe} again
\end{itemize}

\begin{tcolorbox}[breakable,colback=teal!4,colframe=teal!35!black,title={Before you move on}]
What does \mr{गुंडाळणे} mean? (\textquotedblleft{}To roll up; to wrap\textquotedblright{}.) Say it once more.
\end{tcolorbox}
